\section{Computer-Aided Proofs}

This section will give a brief excursion into computer-aided cryptography and automated provers.
Beforehand, it is important to look at different adversarial models in cryptography, as different proof tools assume different types of adversaries.

\subsection{Computationally Unbounded Adversaries}
Computationally unbounded adversaries are unrestricted in their runtime, memory and computational power.
That is, schemes that are secure in the presence of such adversaries are \emph{information-theoretically secure}
and security proofs rely on statistical reasoning.

One scheme secure against computationally unbounded adversaries is the \emph{One-time Pad}~(OTP).
OTP is a symmetric encryption scheme, where messages, ciphertexts and keys have the same length.
Encryption and decryption is then just the XOR of the key and the message or ciphertext.
See Figure~\ref{fig:otp} for the pseudocode.

OTP is information-theoretically secure in the sense that, if the key is truly random and used only to encrypt one message,
no information about the message will be leaked by the ciphertext.
Intuitively this is true, because the ciphertext is just the XOR of the message and a random bit-string (the key).

\begin{figure}[!ht]
    \centering
    \nicoresetlinenr%
    \fbox{%
        \scalebox{\codescalefactor}{%
            %!TEX root=../main.tex
\markersetlen{ndL}{100pt}%
\markersetlen{ndM}{100pt}%
\markersetlen{ndR}{100pt}%
\begin{tabular}[t]{lll}
    \nicodemusbox{\markerlenndL}{%
        \textbf{Proc} $\OTP.\gen$
        \begin{nicodemus}
            \item $k\getsr\{0,1\}^l$
            \item Return $k$
        \end{nicodemus}%
    }%
    &\nicodemusbox{\markerlenndM}{%
        \textbf{Proc} $\OTP.\enc(k, m)$
        \begin{nicodemus}
            \item $c\gets k\oplus m$
            \item Return $c$
        \end{nicodemus}%
    }%
    &\nicodemusbox{\markerlenndR}{%
        \textbf{Proc} $\OTP.\dec(k, c)$
        \begin{nicodemus}
            \item $m\gets k\oplus c$
            \item Return $m$
        \end{nicodemus}%
    }%
\end{tabular}%%
        }%
    }
    \caption{%
        The One-time Pad encryption scheme. The $\oplus$ operator denotes a bitwise \emph{XOR} operation.
    }
    \label{fig:otp}
\end{figure}

There also exist some MAC schemes that are secure in the presence of unbounded adversaries (\alert{Add Sources}).
Statistical events, like collisions of random numbers, are also information-theoretically secure.

Most cryptographic primitives however, are only secure in the presence of computationally bounded adversaries.

\subsection{Computationally Bounded Adversaries}
In contrast to computationally unbounded adversaries, bounded adversaries are restricted in their runtime and memory.
In asymptotic cryptography, we speak of adversaries that run in polynomial time in the security parameter.
These adversaries are still unbounded in their computation power, in the sense that they are as powerful as Turing machines.


